% Original preservation proof for the GD family, 2026-11-14 (Pacific time).
% Requires amsmath, amssymb, booktabs; no private macros are used.
% Numerical-calculus and arithmetic-template definitions are attributed to
% the pinned ProveIt PA194 packet and the updated AX11H calibration packet.
\section{The GD family: independent diagonal and binder parameters}
\label{sec:gd-preservation}

The new source combines explicit-binder generalization with an
accumulated-parameter diagonal axiom.  Its generalization entry is
$U(c,h)$, while its diagonal entry is $\Gamma(d+x)$; the two parameters
are independently programmable.  This section proves preservation for
the sum-test source GD and the two cached-test sources GDH and GDHS.
Their concrete tables and finite binding certificates are supplied
separately.  The proof uses the same numerical calculus as the primary
PA194 machine, so the updated AX11H written absolute calibration can be
reused without changing its PA proof predicate.

\subsection{The exact ten-entry, three-operand interface}

Use the Cantor pairing function and numerical constructors
\[
 p(a,b)=\frac{(a+b)(a+b+1)}2+b,\qquad
 I(A,B)=4p(A,B),\quad M(a,b)=4p(a,b)+1,
\]
\[
 U(i,A)=4p(A,i)+2,\qquad E(a,b)=4p(a,b)+3.
\]
Truth is zero and the refutation target is one.  Set
$\neg A=I(A,1)$, $A\wedge B=\neg I(A,\neg B)$, and
$\exists_i A=\neg U(i,\neg A)$.  The inherited calculus
$\mathcal C_4$ consists of truth, numerical modus ponens, theorem
generalization $A\mapsto U(i,A)$, and these schemes:
\begin{align*}
 h_1(D,H)&=I(\exists_D U(D,H),H),\\
 h_2(D,c,H)&=I(U(D,U(c,H)),U(c,U(D,H))),\\
 h_3(D,c,H)&=I(M(H,D),I(E(H,c),M(c,D))),\\
 h_4(D,c,H)&=I(E(H,D),I(M(c,H),M(c,D))),\\
 h_5(D,H)&=\exists_D E(H,D),\\
 h_6(D,c,H)&=I(U(D,c),I(U(D,I(c,H)),U(D,H))),\\
 \rho(D,c,H,u)&=I(I(I(H,u),D),I(I(D,H),I(c,H))),
       \qquad u\le D,\\
 \Gamma(B)&=I(B,U(B,B)).
\end{align*}
Here $h_6$ is the ordinary distribution head, not its exchanged-premise
variant.  For the guarded scheme put
\begin{align*}
 P(a,b,c)&=I(c,M(c,b)),&
 V(a,b,c)&=I(M(a,b),U(c,M(a,b))),\\
 S(a,b,c)&=U\bigl(b,I(M(a,b),P(a,b,c))\wedge I(M(c,b),c)\bigr),\\
 Q(a,b,c)&=U(c,I(S(a,b,c),\neg P(a,b,c))),\\
 g(a,b,c)&=
 \begin{cases}
 \neg I(V(a,b,c),Q(a,b,c)),&c\ne a\text{ and }c\ne b,\\
 0,&\text{otherwise.}
 \end{cases}
\end{align*}
All these are schemes on exact numerical codes.  In particular,
$I(0,0)=0$ is an actual alias, not an omitted exceptional case.

Write $[a;w]=p(a,w)$ and extend this notation by nesting to the right.
The total projections satisfy
$\operatorname{hd}(0)=\operatorname{tl}(0)=0$ and the usual projection
equations, including $[0;0]=0$.  A store is $(t,w,d,h)$, comprising a
visible top, protected numerical stack, offset, and cached formula.
For the three operands $c,x,y$, set $D=d+x$ and $H=h+y$ and define
\[
\begin{array}{c|l}
q_0&\Gamma(D)\\
q_1&h_6(D,c,H)\\
q_2&h_5(D,H)\\
q_3&h_4(D,c,H)\\
q_4&h_3(D,c,H)\\
q_5&h_2(D,c,H)\\
q_6&h_1(D,H)\\
q_7&\rho(D,c,H,d)\\
q_8&U(c,h).
\end{array}
\]
The exact GD commands are
\begin{align*}
 \mathsf A_{c,x,y}(t,w,d,h)
   &=(g(H,D,c),[q_0,q_1,\ldots,q_8,t;w],0,0),
                                                        \tag{GD-A}\label{eq:gd-ax}\\
 \mathsf M_x(t,w,d,h)
   &=\begin{cases}
       (d+x,w,d+x,0),&I(t,d+x)=h,\\
       (0,w,d+x,0),&I(t,d+x)\ne h.
     \end{cases}                                       \tag{GD-M}\label{eq:gd-mp}
\end{align*}
GD uses
\[
 \mathsf S_\Sigma(t,w,d,h)=
 \begin{cases}
 \mathsf{halt},&t+h=1,\\
 (h,\operatorname{tl}(w),d,\operatorname{hd}(w)),&t+h\ne1,
 \end{cases}                                            \tag{GD-S}\label{eq:gd-sum}
\]
and GDH and GDHS both implement
\[
 \mathsf S_H(t,w,d,h)=
 \begin{cases}
 \mathsf{halt},&h=1,\\
 (h,\operatorname{tl}(w),d,\operatorname{hd}(w)),&h\ne1.
 \end{cases}                                            \tag{GD-H}\label{eq:gd-cache}
\]
In the remainder, $\mathsf S$ means the chosen test consistently.

There are exactly ten newly pushed entries in (GD-A): nine $q$ entries
and the saved old top.  The guarded head remains visible.  The source
retains its closing break and does not fall through the full MP reader.
Consequently AX has exactly the three operands displayed above, its
outgoing offset is zero, and its residual program is the continuation
after the $y,x$ fields.  No fourth operand is consumed.  This distinction
changes both the word encoding and every stack-word exponent from the
eleven-entry AX11 families.

\subsection{Binding the new emitter and the two tests}

The changed emitter performs the same finite constructors and pushes as
its parent.  The first saved universal uses the command parameter $c$
instead of the incoming offset $d$.  Three later parameter copies use
$D$ instead of $c$ to construct $I(D,U(D,D))$.  The other eight emitted
heads, the guard comparisons, and the outgoing cleared scratch registers
are unchanged.  Independent free-term execution checks the complete
output tuple, rather than testing the formula only at sample integers.
Its precondition and endpoint, in register order $(N,T,W,A,D,H,X,Y,L)$,
are
\[
 (n,t,w,c,d,h,0,0,p(y,p(x,\ell)))
 \longmapsto
 (n,g(H,D,c),[q_0,\ldots,q_8,t;w],0,0,0,0,0,\ell).
                                                        \tag{GD-End}\label{eq:gd-end}
\]
The original breaks remain literal breaks.  In particular this proof
does not use a historical-zero invariant to replace a break by a
conditional decrement.

The independent \texttt{check\_source.py} generates every primitive
instruction and outgoing edge from handwritten arithmetic templates;
it imports neither the construction compiler nor a Turing-machine
simulator.  Its arrays equal the delivered 3,654-site GD, 3,653-site
GDH, and 3,650-site GDHS sources.  As a separate check it expands the
named binary source tree and resolves all jumps by subtree intervals,
obtaining those same arrays.  The respective template decompositions
contain 942, 943, and 941 instances.  The common command ports are
fetch 25, decode 34, AX 50, full MP 3554, MP comparison 3624,
successful MP copy 3634, return 3641, and Shift 3642.
The AX closing break at 3553 jumps to the shared return at 3641,
which jumps to 25; its command clock includes both steps.

The same checker executes AX in a free algebra with all input
parameters symbolic.  It checks both guard comparisons, every
intermediate scratch precondition, and all nine final registers in
(GD-End).  For each source it also replays 882 local exact-duration
examples and 324 composed Shift/MP examples, covering both comparison
outcomes and both fresh and positive-continuation fetch cases.
These examples check the arithmetic implementations; the uniform
counter-induction proof is the reason the clock equations hold for
arbitrary natural inputs.

The H source clears $T$, copies and restores $H$ into $T$, and increments
the initially zero $X$.  Its root comparison therefore sees
$(T,H,X)=(h,h,1)$.  The HS source instead duplicates $H$ into $T,X$ and
increments $H$, giving $(T,H,X)=(h,1,h)$.  The inherited destructive
comparison theorem gives equality precisely at $h=1$ in both cases;
on inequality it clears the compared registers and reaches the same
stack-pop continuation.  Both sources therefore implement (GD-H) for
every store, despite having different instruction arrays and clocks.

For clarity, if $C(a,b)$ denotes the exact inherited comparison clock,
the respective durations before the common stack-pop or halt suffix are
\[
 (2t+1)+(7h+2)+2+C(h,1),\qquad
 (2t+1)+(4h+1)+2+C(1,h).
\]
Here clear, transfer, duplicate, and copy-and-restore cost
$2r+1$, $3r+1$, $4r+1$, and $7r+2$ respectively.  For the sum test,
the unchanged parent source transfers $T$ into $X$, duplicates $H$ into
$T,X$, then transfers $X$ into $H$.  Thus it reaches
$(T,H,X)=(h,t+h,0)$ before incrementing $X$ and jumping to the root.
It compares the sum with one while retaining the old cache as the next
top, and its corresponding prefix clock is
\[
          (3t+1)+(4h+1)+(3(t+h)+1)+2+C(t+h,1).
\]
These expressions are composed with the pairing, unpairing, and fetch
clocks already proved in the inherited source-command template argument.
That argument uses quantifier-free induction on completed loop rounds,
with register equalities and partial-sum clocks as its invariant.  The
new source changes only a finite composition of those templates.  Thus
the same proof applies once the delivered instruction and edge arrays
are checked against that composition.

\subsection{Stack-word equations, including the zero alias}

Every $q_0,\ldots,q_7$ and $g(H,D,c)$ is a multiple of four;
$q_8=U(c,h)$ is congruent to two modulo four and is at least two.
Direct iteration of the command equations gives, without a halt,
\begin{align*}
 \mathsf A_{c,x,y};\mathsf S^9(t,w,d,h)
    &=(q_7,[t;w],0,q_8),\\
 \mathsf A_{c,x,y};\mathsf S^{10}(t,w,d,h)
    &=(q_8,w,0,t).
\end{align*}
All guards tested in these two calculations concern the safe heads and
zero; neither calculation tests the saved old top alone.

Put $A_0=\mathsf A_{0,0,0}$ and define
\[
 \mathsf{Push}=A_0;\mathsf S^9;\mathsf M_0,\qquad
 \mathsf{Pop}=A_0;\mathsf S^{11},\qquad
 \mathsf{Discard}=\mathsf M_0.
\]
For either Shift variant,
\begin{align*}
 \mathsf{Push}(t,w,d,h)&=(0,[t;w],0,0),\\
 \mathsf{Discard}(t,w,0,h)&=(0,w,0,0),\\
 \mathsf M_i(0,w,0,0)&=(0,w,i,0).
\end{align*}
Push's last MP compares an implication code of residue zero with $q_8$
of residue two, and therefore fails.  Discard has result zero whichever
comparison branch is taken.  The last identity succeeds with result zero
at $i=0$ and fails at $i>0$, since $I(0,i)>0$ then.

The Pop laws differ between the two tests:
\begin{align*}
 \mathsf{Pop}_\Sigma(t,[a;w],d,h)&=(t,w,0,a),\\
 \mathsf{Pop}_H(t,[a;w],d,h)&=
   \begin{cases}\mathsf{halt},&t=1,\\(t,w,0,a),&t\ne1.
   \end{cases}                                         \tag{GD-Pop}\label{eq:gd-pop}
\end{align*}
The eleventh Shift sees top $q_8$ and cache $t$.  Its sum is at least two,
so the sum test cannot halt.  The cache test halts exactly at $t=1$.
These equations are equations for the numerical pairing and total
projections.  They hold when one or several apparent stack cells are
zero, including $[0;0]=0$; no free-list assumption is made.

\subsection{Selecting every inherited axiom instance}

For $0\le j\le7$, first take incoming offset and cache zero.  Apart
from the fourth parameter of $\rho$, which is treated below, the word
\[
 \mathsf{Discard};\mathsf A_{c,D,H};
             \mathsf S^{j+2};\mathsf{Pop}^{8-j}           \tag{GD-Sel}\label{eq:gd-sel}
\]
returns $(q_j,w,0,0)$ from every $(t,w,0,0)$.  The first Shifts put
$q_j$ at the top.  Every Pop preserves that top and removes one remaining
generated stack entry; the last such entry is the saved zero.
Since $q_j\ne1$, no cached-test Pop halts.  For the guarded head use
\[
 \mathsf{Discard};\mathsf A_{c,D,H};\mathsf{Pop}^{10}.
\]
It returns $(g(H,D,c),w,0,0)$, again without a halt.  A failed guard
returns truth, as in the parent machine.

For the desired dominated instance $\rho(D,c,H,u)$ with $u\le D$, use
\[
 \mathsf{Discard};\mathsf M_u;
 \mathsf A_{c,D-u,H};\mathsf S^9;\mathsf{Pop}.
\]
The actual accumulated parameter is $u+(D-u)=D$ and the incoming offset
recorded by the fourth argument of $\rho$ is exactly $u$.

The free diagonal parameter requires a deliberate adjustment.
The new diagonal head is $\Gamma(D)$, so for an arbitrary prescribed
$B$ its selector is
\[
 \mathsf{Discard};\mathsf A_{0,B,0};
                        \mathsf S^2;\mathsf{Pop}^8.      \tag{GD-Diag}\label{eq:gd-diag}
\]
This sets $D=B$ directly.  Thus the source still emits every instance
$\Gamma(B)$ of the same scheme; no diagonal axiom has been lost.
The non-diagonal selectors set $c,D,H$ as before, independently.

Explicit-binder generalization is implemented by
\[
 \mathsf{Gen}_i=
       \mathsf{Push};\mathsf S;\mathsf A_{i,0,0};\mathsf S^{10}.
                                                        \tag{GD-Gen}\label{eq:gd-gen}
\]
From $(a,w,0,0)$, Push and the next Shift give $(0,w,0,a)$.
That Shift tests zero before moving $a$ into the cache.  AX saves $U(i,a)$
and clears the visible cache; the final ten Shifts return
$(U(i,a),w,0,0)$.  This computation is nonhalting even if $a=1$.
The changed diagonal parameter does not affect it.

\subsection{A compiler that tolerates early refutations}

Suppose that a major premise proves $I(a,b)$ and a minor premise proves
$a$.  On the store containing the two premises and a reserved zero,
consider
\[
 \mathsf S;\mathsf S;\mathsf M_b;\mathsf{Pop}
       \quad\hbox{from}\quad(0,[a,I(a,b),0;w],0,0).
                                                        \tag{GD-MPs}\label{eq:gd-mps}
\]
The first Shift is safe.  The second halts if $a=1$ for either test.
Otherwise it leaves $(a,[0;w],0,I(a,b))$ and MP succeeds, giving
$(b,[0;w],b,0)$.  Sum-test Pop then returns $(b,w,0,0)$ unconditionally;
cached-test Pop halts at $b=1$ and otherwise returns the same store.
Consequently (GD-MPs) always either halts or returns $(b,w,0,0)$.
An early halt is an allowed successful outcome of the refutation
compiler, rather than a case excluded by assuming consistency.

Represent numerical derivations by finite lists with earlier premise
indices.  Use Discard for truth, the preceding selectors for axiom
instances, and append $\mathsf{Gen}_i$ at a generalization line.  At a
modus-ponens line with major derivation $e_1$ and minor derivation $e_0$
define the finite word
\[
\begin{split}
 W_e={}&\mathsf{Discard};\mathsf{Push};W_{e_1};\mathsf{Push};
 W_{e_0};\mathsf{Push};\mathsf S;\mathsf S;
 \mathsf M_b;\mathsf{Pop}.
\end{split}                                             \tag{GD-CmpDef}\label{eq:gd-cmpdef}
\]
The first two commands reserve the zero below the premises.  Recursive
words preserve the protected stack when they return.  The local laws
therefore establish
\[
 W_e(t,w,0,0)\ \text{halts or returns}\
             (\operatorname{concl}(e),w,0,0)             \tag{GD-Cmp}\label{eq:gd-cmp}
\]
for every valid numerical derivation and every $t,w$.

A common finishing word works for both tests:
\[
                    \mathsf{Finish}=\mathsf{Push};\mathsf S;\mathsf S.
\]
On $(1,w,0,0)$, Push returns $(0,[1;w],0,0)$, the first Shift gives
$(0,w,0,1)$, and the next Shift halts.  Thus for every refutation $e$,
\[
                   V_e=A_0;W_e;\mathsf{Finish}            \tag{GD-Ref}\label{eq:gd-ref}
\]
halts from every incoming store.  The initial AX makes offset and cache
zero; its otherwise arbitrary output and protected stack are permitted
inputs to (GD-Cmp).

\subsection{Why the compiler proof uses only allowed induction}

An unbounded assertion about all stores is not used as an induction
formula.  Fix a valid list of $N>0$ proof lines and compile its words
$W_0,\ldots,W_{N-1}$ bottom-up by primitive recursion.  Let $L\ge1$
bound their lengths and $K$ bound all their command codes.  Encode stores
and an absorbing halt by natural numbers.  For the primitive recursive
command step $\operatorname{next}$ define
\[
 \Phi_K(s)=1+\max\bigl(s,
            \max_{c\le K,\ u\le s}\operatorname{next}(c,u)\bigr).
\]
This is monotone, and bounded maxima and iteration are primitive
recursive.  Fix the desired normalized input code $s_0$ and put
\[
              B(j)=\Phi_K^{\,L(N-1-j)}(s_0)\qquad(j<N).
\]
An occurrence of a premise word inside $W_j$, started at input at most
$B(j)$, receives an input at most $\Phi_K^L(B(j))$; all its preceding
commands belong to $W_j$.  If its premise index is $i<j$, this is at most
$B(i)$.

At induction stage $j$, take the bounded conjunction of (GD-Cmp) for
every earlier line $i<j$ and every normalized input code at most $B(i)$.
Finite word execution, validity, normalization, and the halt flag are
primitive recursive.  Hence the whole bounded conjunction is a single
primitive recursive Boolean equation.  Its quantifier-free induction
step follows from the preceding local word equations and the displayed
input bound.  At the last line, $B(N-1)=s_0$.  Generalizing the arbitrary
input after this induction yields (GD-Cmp), and hence (GD-Ref), in iPRA.
All word construction and execution witnesses are explicit primitive
recursive functions.  No least successful execution time is used.

\subsection{Three-operand encoding and the carried-store scheduler}

The correct word encoder for this source is
\[
 e(\epsilon)=0,\qquad e(a::v)=1+b_a(e(v)),
\]
\[
 b_{\mathsf S}(r)=p(0,r),\qquad
 b_{\mathsf M_x}(r)=p(1,p(x,r)),\qquad
 b_{\mathsf A_{c,x,y}}(r)=p(c+2,p(y,p(x,r))).              \tag{GD-Enc}\label{eq:gd-enc}
\]
There is no fourth nested AX field.  The exact live scheduler state is
$(n,\ell,t,w,d,h)$; its initial value is $(2,0,0,0,0,0)$.  If $\ell=0$
the next raw seed is $n$ and the new fresh counter is $n+1$.  If $\ell>0$
the next raw seed is $\ell-1$ and the fresh counter stays $n$.  After
executing one decoded command, the residual after its fields is the new
continuation.  The store is carried across all candidate boundaries.

Every residual is a nested Cantor component of the raw seed and is at
most that seed.  Thus a positive continuation strictly decreases, and a
candidate of raw seed $n$ drains within $n+1$ command steps unless it
halts.  These assertions follow by quantifier-free induction on bounded
command iterations.

Prefix $\mathsf A_{k,0,0}$ to the all-store halting suffix $V_e$.
The raw seed of the resulting candidate is
\[
                  N_k=p(k+2,p(0,p(0,e(V_e)))).            \tag{GD-Seed}\label{eq:gd-seed}
\]
It is at least $k+2$ and at least three.  These seeds are cofinal in the
fresh enumeration, and the arbitrary prefix output causes no difficulty
because $V_e$ halts from every store.  From the clean start choose $k=0$;
draining the candidates through $N_0$ takes at most $(N_0+2)^2$ abstract
steps unless a halt has occurred earlier.  This is an explicit primitive
recursive witness bound.  It never assumes that failed trials reset the
store.  Therefore iPRA proves that a $\mathcal C_4$ refutation yields a
finite scheduler halt.

\subsection{Soundness and transfer of the absolute calibration}

For soundness, annotate the visible top, cache, and every canonical
numerical stack entry by indices in a finite valid $\mathcal C_4$
derivation.  Truth supplies the initial annotations and all virtual
zero entries.  For $w>0$, $\operatorname{tl}(w)<w$, so canonical stack
annotation and its validity are bounded primitive recursive operations.

AX appends valid proofs of its entries.  The new $\Gamma(D)$ is an
instance of the unchanged diagonal scheme.  Its $U(c,h)$ follows by
theorem generalization from the annotated cache.  The dominated instance
has $d\le d+x$; the other heads are their displayed schemes, and a failed
gate supplies truth.  The saved top was already proved.  Successful MP
uses the proofs of $t$ and $I(t,d+x)$; failure supplies truth.  A
nonhalting Shift merely moves annotated formulas.  A cached-test halt
exposes the proof of $h=1$.  A sum-test halt has $t+h=1$, so either
$t=1,h=0$ or $t=0,h=1$ by the decidable natural-number case analysis;
one of the two annotations is therefore a refutation.

The annotation update is primitive recursive and its validity checker
is Boolean.  Quantifier-free induction on a finite command run extracts
a numerical refutation from every scheduler halt.  Combined with the
compiler, for each of the two abstract tests this proves
\[
 \mathrm{iPRA}\vdash\quad
 \operatorname{Halt}(R_{GD})\ \longleftrightarrow\
           \exists e\,\operatorname{Prf}_{\mathcal C_4}(e,1).
                                                        \tag{GD-Calc}\label{eq:gd-calc}
\]

The source-command clocks are total primitive recursive functions.
Composing their positive live-command segments gives a cumulative
source clock.  The inherited physical proof supplies total primitive
recursive macro clocks, including frame, gap, variable file width, and
spare-file growth.  Instantiating it with the delivered decoder, folds,
source instructions, and complete blank-start certificate gives the
physical/source equivalence.  Halt is absorbing; its subsequent clock
increments may be zero.  Conversely, a proposed first physical or
source halt cannot occur within more positive nonhalting segments than
its own time.  This is the same conditional cofinality argument used in
the inherited physical audit, and does not assume positive durations
after halt.

Let $M$ be one of the concrete GD, GDH, or GDHS tables with those
bindings checked.  Composition with (GD-Calc) gives
\[
 \mathrm{iPRA}\vdash\quad
 \operatorname{Halt}_{\rm blank}(M)\ \longleftrightarrow\
           \exists e\,\operatorname{Prf}_{\mathcal C_4}(e,1).
                                                        \tag{GD-Phys}\label{eq:gd-phys}
\]
The numerical calculus in this formula is exactly the ordinary
$\mathcal C_4$ calibrated by the updated AX11H packet.  In particular
the new selector (GD-Diag) ranges over the same diagonal scheme, and no
logical axiom or inference rule has been strengthened or weakened.

The inherited absolute calibration uses its Foundation-to-ordinary-PA
adapter, finite-demand logical compiler, model-uniform numerical
soundness, and the pinned one-way negative-transfer endpoints.  It
establishes the negative calibration
\[
 \mathrm{iPRA}\vdash\quad
 \operatorname{Con}(\mathrm{PA})\ \longleftrightarrow\
       \neg\exists e\,\operatorname{Prf}_{\mathcal C_4}(e,1).
                                                        \tag{GD-Cal}\label{eq:gd-cal}
\]
Equivalently, (GD-Cal) follows by combining its written AX11H
consistency/nonhalting theorem with that source's proved equivalence to
the same numerical refutation predicate.  This use of the inherited
calibration does not assert an intuitionistic equivalence between two
positive existential PA and numerical proof predicates.  In particular
it does not extract a positive witness from a double negation.

Applying negation to the two implications of (GD-Phys), then using
(GD-Cal), is constructive and gives the claimed preservation theorem:
\[
 \boxed{\quad\mathrm{iPRA}\vdash\quad
   \operatorname{Con}(\mathrm{PA})\ \longleftrightarrow\
                 \neg\operatorname{Halt}_{\rm blank}(M).\quad}
\]
No consistency assumption is used to prove this equivalence.  As with
the inherited absolute calibration, this is a written mathematical
proof whose finite table and source bindings have executable checks;
it is not a completed end-to-end Lean kernel certification.

\subsection{Why GD is a different source architecture}

In the earlier explicit-binder G source the same command parameter
supplied both $U(c,h)$ and $\Gamma(c)$.  GD retains the binder $c$ but
uses the independent accumulated value $d+x$ in the diagonal entry.
For example, from offset and cache zero, take $c=0,x=1,y=0$.
Both sources save $U(0,0)=2$, while the older G source emits
$\Gamma(0)=20$ and GD emits $\Gamma(1)=832$.
These are different command outputs on the same input.  A uniform
renaming that swapped $c$ and $D$ would also change the generalization
entry to $U(D,h)$, which GD does not do.  Its independence of diagonal
and binder parameters is therefore more than a global parameter rename.
The three delivered tables additionally use distinct literal Shift
implementations, with GDH and GDHS sharing an abstract test and GD
having the different sum-test transition relation.
